We define the vector of decision variables as:
\[
X = \begin{pmatrix}
x_{11} & x_{12} & x_{13} & x_{21} & x_{22} & x_{23} & x_{31} & x_{32} & x_{33} &
y_1 & y_2 & y_3 & r_1 & r_2 & r_3 & c_1 & c_2 & c_3
\end{pmatrix}^{\!\top}.
\]
where \(x_{ij}\) represents the quantity shipped from distribution center
\(\mathrm{DC}_i\) to city \(j\). We assume that \(x_{i1}\) corresponds to Grenoble, \(x_{i2}\) to Lyon, and
\(x_{i3}\) to Valence.

\subsubsection{Objective function}

The objective is to minimize the total cost, which includes transportation,
labor, and inventory costs, as well as penalties related to carbon emissions
and renewable energy use.

The transportation cost is:
\[
C_{\mathrm{transport}} =
4x_{11} + 5x_{12} + 6x_{13}
+ 7x_{21} + 8x_{22} + 9x_{23}
+ 10x_{31} + 11x_{32} + 12x_{33}.
\]

The labor cost is:
\[
C_{\mathrm{labor}} = 20y_1 + 22y_2 + 25y_3.
\]

The inventory cost is:
\[
C_{\mathrm{inventory}} =
0.10\,(x_{11} + x_{12} + x_{13})
+ 0.20\,(x_{21} + x_{22} + x_{23})
+ 0.30\,(x_{31} + x_{32} + x_{33}).
\]

The penalty associated with carbon emissions is:
\[
C_{\mathrm{carbon}} = 50\,(c_1 + c_2 + c_3).
\]

Finally, the cost associated with renewable energy is:
\[
C_{\mathrm{renewable}} = 10\,(r_1 + r_2 + r_3).
\]

The objective function is therefore written as:
\[
\boxed{
\begin{aligned}
\min \quad Z ={}& C_{\mathrm{transport}} + C_{\mathrm{labor}} + C_{\mathrm{inventory}}
+ C_{\mathrm{carbon}} + C_{\mathrm{renewable}} \\
={}& 4x_{11} + 5x_{12} + 6x_{13} + 7x_{21} + 8x_{22} + 9x_{23}
+ 10x_{31} + 11x_{32} + 12x_{33} \\
&+ 20y_1 + 22y_2 + 25y_3 \\
&+ 0.10\,(x_{11} + x_{12} + x_{13})
+ 0.20\,(x_{21} + x_{22} + x_{23})
+ 0.30\,(x_{31} + x_{32} + x_{33}) \\
&+ 50\,(c_1 + c_2 + c_3) + 10\,(r_1 + r_2 + r_3).
\end{aligned}
}
\]

The objective function can be written as:\(Z(X) = f^{\top} X\), avec
\[
f = \begin{pmatrix}
\underbrace{4 \;\; 5 \;\; 6 \;\; 7 \;\; 8 \;\; 9 \;\; 10 \;\; 11 \;\; 12}_{x}
\;\; \underbrace{20 \;\; 22 \;\; 25}_{y}
\;\; \underbrace{10 \;\; 10 \;\; 10}_{r}
\;\; \underbrace{50 \;\; 50 \;\; 50}_{c}
\end{pmatrix}^{\!\top}.
\]


\subsubsection{Capacity constraints}

The maximum capacities of the three distribution centers are
\(2000\), \(2500\), and \(3000\) units, respectively. The constraints
are therefore written as:
\[
\begin{aligned}
x_{11} + x_{12} + x_{13} &\leq 2000, \\
x_{21} + x_{22} + x_{23} &\leq 2500, \\
x_{31} + x_{32} + x_{33} &\leq 3000.
\end{aligned}
\]

In matrix form, these constraints are written as:
\[
A_{\mathrm{cap}}\, X \leq b_{\mathrm{cap}},
\]
with
\[
A_{\mathrm{cap}} =
\begin{bmatrix}
1 & 1 & 1 & 0 & 0 & 0 & 0 & 0 & 0 & 0 & \cdots & 0 \\
0 & 0 & 0 & 1 & 1 & 1 & 0 & 0 & 0 & 0 & \cdots & 0 \\
0 & 0 & 0 & 0 & 0 & 0 & 1 & 1 & 1 & 0 & \cdots & 0
\end{bmatrix}
 and
\quad
b_{\mathrm{cap}} =
\begin{bmatrix}
2000 \\ 2500 \\ 3000
\end{bmatrix}.
\]

\subsubsection{Demand satisfaction constraints}

The demands of the three cities are expressed by:
\[
\begin{aligned}
x_{11} + x_{21} + x_{31} &= 1500, \\
x_{12} + x_{22} + x_{32} &= 1000, \\
x_{13} + x_{23} + x_{33} &= 1200.
\end{aligned}
\]

In matrix form, these constraints are written as:
\[
A_{\mathrm{dem}}\, X = b_{\mathrm{dem}},
\]
with
\[
A_{\mathrm{dem}}  =
\begin{bmatrix}
1 & 0 & 0 & 1 & 0 & 0 & 1 & 0 & 0 & 0 & \cdots & 0 \\
0 & 1 & 0 & 0 & 1 & 0 & 0 & 1 & 0 & 0 & \cdots & 0 \\
0 & 0 & 1 & 0 & 0 & 1 & 0 & 0 & 1 & 0 & \cdots & 0
\end{bmatrix}
\quad \text{and} \quad
b_{\mathrm{dem}} =
\begin{bmatrix}
1500 \\ 1000 \\ 1200
\end{bmatrix}.
\]

\subsubsection{Working hours constraints}

The working hour limits are:
\[
y_1 \leq 45, \qquad y_2 \leq 55, \qquad y_3 \leq 50,
\]

That s mean :
\[
A_{\mathrm{lab}}\, X \leq b_{\mathrm{lab}},
\]
with
A_{\mathrm{lab}} =
\begin{bmatrix}
0 & \cdots & 0 & 1 & 0 & 0 & 0 & \cdots & 0 \\
0 & \cdots & 0 & 0 & 1 & 0 & 0 & \cdots & 0 \\
0 & \cdots & 0 & 0 & 0 & 1 & 0 & \cdots & 0
\end{bmatrix}
\]
and 
\[
b_{\mathrm{lab}} =
\begin{pmatrix}
45 \\ 55 \\ 50
\end{pmatrix}.
\]

\subsubsection{Renewable energy constraints}

The total amount of renewable energy used must be
greater than or equal to \(1000\) kWh:
\[
\boxed{r_1 + r_2 + r_3 \geq 1000,}
\]
In other words:
\[
A_{\mathrm{ren}}\, X = -(r_1 + r_2 + r_3) \leq -1000.
\]
with : 
\[
A_{\mathrm{ren}} =
\begin{bmatrix}
0 & \cdots & 0 & -1 & -1 & -1 & 0 & 0 & 0
\end{bmatrix}
  and
\qquad
b_{\mathrm{ren}} = -1000.
\]
\subsubsection{Carbon emission constraints}

Each unit transported generates \(0.05\) tonnes of carbon.
The emissions of each center are therefore:
\[
\begin{aligned}
c_1 &= 0.05\,(x_{11} + x_{12} + x_{13}), \\
c_2 &= 0.05\,(x_{21} + x_{22} + x_{23}), \\
c_3 &= 0.05\,(x_{31} + x_{32} + x_{33}).
\end{aligned}
\]

The total amount of emissions must not exceed \(300\) tonnes:
\[
\boxed{c_1 + c_2 + c_3 \leq 300,}
\]
so 
\[
A_{\mathrm{CO_2}} =
\begin{pmatrix}
0 & \cdots & 0 & 1 & 1 & 1
\end{pmatrix} 
and
\qquad
b_{\mathrm{CO_2}} = 300.
\]

\subsection{Matrix reformulation of the problem}

By grouping all the previous constraints, the problem can be reformulated
in the standard form of a linear programming problem:
\[
\boxed{
\begin{aligned}
\min_{X} \quad & f^{\top} X \\
\text{subject to} \quad
& A\, X \leq b, \\
& A_{\mathrm{eq}}\, X = b_{\mathrm{eq}}, \\
& X \geq 0,
\end{aligned}
}
\]
where \(X \in \mathbb{R}^{18}\) is the vector of decision variables defined
previously and \(f \in \mathbb{R}^{18}\) is the vector of objective function
coefficients.

The matrix \(A\) gathers the inequality constraints: capacity, working
hours, renewable energy, and carbon emissions. The matrix
\(A_{\mathrm{eq}}\) gathers the equality constraints: demand satisfaction and
the relations defining carbon emissions. We thus have:
\[
A =
\begin{pmatrix}
A_{\mathrm{cap}} \\
A_{\mathrm{lab}} \\
A_{\mathrm{ren}} \\
A_{\mathrm{CO_2}}
\end{pmatrix},
\qquad
b =
\begin{pmatrix}
b_{\mathrm{cap}} \\
b_{\mathrm{lab}} \\
b_{\mathrm{ren}} \\
b_{\mathrm{CO_2}}
\end{pmatrix},
\]
and, similarly,
\[
A_{\mathrm{eq}} =
\begin{pmatrix}
A_{\mathrm{dem}} \\
A_{\mathrm{em}}
\end{pmatrix},
\qquad
b_{\mathrm{eq}} =
\begin{pmatrix}
b_{\mathrm{dem}} \\
0_{3 \times 1}
\end{pmatrix}.
\]

with :
\[
A_{\mathrm{em}} =
\begin{bmatrix}
-0{,}05 & -0{,}05 & -0{,}05 & 0 & 0 & 0 & 0 & 0 & 0 & 0 & \cdots & 0 & 1 & 0 & 0 \\
0 & 0 & 0 & -0{,}05 & -0{,}05 & -0{,}05 & 0 & 0 & 0 & 0 & \cdots & 0 & 0 & 1 & 0 \\
0 & 0 & 0 & 0 & 0 & 0 & -0{,}05 & -0{,}05 & -0{,}05 & 0 & \cdots & 0 & 0 & 0 & 1
\end{bmatrix}
and 
\qquad
b_{\mathrm{em}} =
\begin{pmatrix}
0 \\ 0 \\ 0
\end{pmatrix}
\]
